\section{Conditional identification: complete proofs}
\label{sec:annex-p16-conditional-interventional-identification}

\begin{annexitem}{annex:p16.idc-positive-context}{thm:p16-idc-positive-context}{Support of all conditional intervention contexts on positive laws}
\begin{proof}
Fix an intervention assignment and an assignment of every remaining
observed variable. Their combination is a full assignment \(v\)
with positive observed mass. In the finite latent-sum formula,
replacing the intervened equations removes their indicator factors.
Every summand is nonnegative, so the resulting intervention mass
on the remaining assignment is at least \(p(v)>0\). Every
partial assignment of remaining coordinates has a full extension,
and its marginal therefore contains a positive term. This gives
all claimed conditioning denominators, regardless of whether the
action values were original inputs or came from exchanged evidence.
\end{proof}
\end{annexitem}

\begin{annexitem}{annex:p16.idc-positive-path-gadget}{thm:p16-idc-positive-path-gadget}{Explicit active-path coordinates and positive perturbation}
\begin{proof}
Choose an active simple path from \(w\) to one vertex of
\(Y\). Split it at its internal colliders, of which there
are \(m\). Each of the \(m+1\) remaining segments has no
internal collider and has a unique source from which its edges
point toward its two ends, allowing a source at an endpoint.
Indeed along a collider-free segment the arrow orientations can
change from backward to forward at most once; a change in the
other direction would create a collider. Assign an independent
fair bit to each segment's source and copy it along the two
arms toward the segment's ends. If the source is a latent root,
use its bit directly; if observed, use a private fair input there.

At each collider, output the exclusive-or of the two adjacent
segment bits. Activity supplies a directed route from that collider
to \(C\), allowing length zero. Stop at the first conditioning
vertex on a chosen route, and copy the collider's exclusive-or
along that route. Allocate a separate coordinate to each copy
channel whenever routes intersect each other or the original path.
All coordinates use existing directed parent edges, and every
equation is evaluated in the underlying DAG's topological order.
Sources use private inputs and ignore other arguments. The output
at \(w\) is just the first segment bit \(K\): no activation
route can pass through this sink to reach \(C\), which excludes
\(w\). At the other endpoint retain its last segment bit \(B\).
Give other vertices a constant coordinate if none is otherwise needed.

Condition all coordinates at \(C\) to zero. A conditioned
vertex cannot be an internal noncollider of the original path,
and neither endpoint belongs to \(C\). Thus its nonconstant
coordinates are only collider exclusive-or values or copies of
such values, never an individual segment bit. The conditioning
event imposes exactly the \(m\) adjacent-segment equalities,
possibly with duplicate copies. It has probability \(2^{-m}\)
and leaves the common segment bit fair. Consequently the
unnormalized two-by-two table for \(K,B,C=0\) has two diagonal
entries \(2^{-(m+1)}\) and two off-diagonal zeros. Its determinant
is \(d_0=2^{-2m-2}>0\).

Now add independent private flip noise of probability \(\delta\)
to every observed output coordinate, using a fresh noise input
for each coordinate. Let \(q\geq1\) be their finite number.
Conditional on any latent roots and earlier parents, every vector
output has positive probability when \(0<\delta<1\). Hence
the full observed law is strictly positive. Couple this construction
to the noiseless one using the same segment source bits. Except
when at least one flip occurs, their entire outputs agree. The
probability of a flip is at most \(q\delta\), so each entry
of the relevant unnormalized table changes by at most that amount.
Since all masses are at most one, its determinant changes by at
most \(4q\delta\). Choose a positive rational \(\delta\)
with \(4q\delta<d_0/2\). The new determinant remains nonzero.
Dividing all four entries by the positive conditioning mass scales
the determinant by a positive number, preserving nonzero determinant.

All added noises are private and independent of the source bits.
Formal latent-root forks correspond to allowed bidirected edges,
and may be implemented by independent pair inputs; this constructs
one permitted model, rather than replacing another model's inputs.
The sink \(w\) still has only its one output coordinate and
no used outgoing edge. This proves the claimed finite positive gadget.
\end{proof}
\end{annexitem}

\begin{annexitem}{annex:p16.idc-unused-conditioning-bit}{thm:p16-idc-unused-conditioning-bit}{A passive relay witness and the full-conditional argument}
\begin{proof}
Joint graph-uniform completeness supplies a hedge for \((A,W)\).
Use the positive core/relay construction of
\cref{thm:p16-hedge-four-state-routing}, with its stated coordinate
elimination. Choose its routing forest to stop at the first vertex
of \(W\) on every route. Write \(J\subseteq W\) for
the routing sinks, and \(L_J\) for their relay bit tuple.
Each such bit is unused in every other equation: it is a routing
sink, and core equations ignore all relay coordinates.

Let \(C_0\) be the tuple of all remaining observed coordinates
at vertices of \(W\). These are retained core coordinates.
The construction omits core coordinates at core-forest roots;
every retained core in \(W\) is therefore a nonroot forest
coordinate or an independent outside coordinate. None depends
on the private noise \(E_0\) used at the selected forest root,
or on any relay noise. Under the action \(A=a\), model 2
has the exact parity identity
\[
 \bigoplus_{j\in J}L_j=E_0\oplus\bigoplus_{i\in H}E_i,
\]
whose expected sign is a constant \(\rho>0\), also conditional
on any \(C_0\). All noises have one-probabilities below one half
and are independent of \(C_0\). In model 1 that relay parity
is fair unconditionally, by the same construction. If the conditional
distributions of \(L_J\) given \(C_0\) agreed everywhere,
model 1 would also have conditional expected sign \(\rho\)
everywhere and hence unconditional expected sign \(\rho\), a
contradiction. Both current laws are positive, so some common
\(C_0=c_0\) gives two different strictly positive distributions
of the relay tuple.

Two positive distributions on a finite binary product that have
the same one-coordinate conditional distributions given all other
coordinates must be equal. To prove this, take their positive
mass ratio. Equality of a coordinate conditional makes this ratio
unchanged when just that coordinate is changed. Any two binary
tuples can be connected by finitely many such changes, so the
ratio is constant; normalization makes it one. Applying the
contrapositive to the two relay distributions gives an index
\(w\in J\) and values of the other relay bits for which
the conditional probability of \(L_w=1\) differs. Together
with \(c_0\), these are exactly all the other evidence coordinates.
The selected bit has the required unused-output property. Equal
positive observations, unchanged declared graph, and positive
post-action contexts are properties already proved for the
core/relay construction and its coordinate projection.
\end{proof}
\end{annexitem}

\begin{annexitem}{annex:p16.idc-terminal-reduction}{thm:p16-idc-terminal-reduction}{The complete two-case reduction with an independent masking gadget}
\begin{proof}
The joint-to-conditional direction is the successful normalization
proof. For necessity, suppose the terminal joint is not identified.
If its conditioning marginal is identified, apply
\cref{thm:p16-idc-known-denominator}. Otherwise use the unused-
conditioning-bit lemma to obtain two positive finite old models,
the same observed law, and a bit \(L\) at \(w\in W\)
whose conditional one-probabilities \(t_1,t_2\) given the same
remaining old evidence tuple \(R=r\) differ. Positivity makes
both \(t_i\) lie strictly between zero and one.

No observation is eligible at the terminal query. Thus \(w\)
is d-connected to some response vertex in the canonical DAG
after cutting incoming arrows at \(A\), outgoing arrows at
\(w\), and removing the fixed action vertices, given
\(W\setminus\{w\}\). Here \(w\) is a sink. Apply the
positive-path gadget lemma in that DAG. It supplies a bit \(K\)
at \(w\), a response bit \(B\), and an assignment \(c\)
of the other conditioning gadget coordinates such that
\[
 h_{kb}=\mathbb P(K=k,B=b,C=c)
 \quad(k,b\in\{0,1\})
\]
is a strictly positive table with nonzero determinant. Give action
vertices independent dummy gadget coordinates ignored by all other
gadget equations. The gadget law on the remaining vertices is
then unchanged by any action values at \(A\).

Take this gadget independently of each old model, using new latent
inputs. At every observed vertex combine the old and gadget output
coordinates, except at \(w\): replace its old bit \(L\)
by \(T=L\oplus K\), and omit \(K\) itself from the
observed output. Every other old coordinate at \(w\) is retained.
This is a legal total structural model on the original declared
graph. The old bit \(L\) is unused by other old equations,
and the gadget bit \(K\) has no used child. Both can be
computed inside the equation at \(w\) and replaced by their
exclusive-or without making an unavailable value a parent elsewhere.
The old and gadget parent edges are allowed edges of the graph;
their independent latent groups introduce no extra bidirected edge.
All alphabets remain finite.

Before this last replacement, the two full observational laws agree,
being the same old law times the same independent gadget law.
Applying the same coordinate map preserves that agreement. The
new law is strictly positive: every proposed combined output has
an extension with, for example, \(K=0,L=T\), and both old
and gadget assignments have positive mass. The same structural
coordinate map applies after \(A=a\), since \(w\notin A\)
and neither omitted coordinate is used by another equation.

Condition in each new model on \(R=r,C=c,T=0\). This is
positive in both models. Independence of the old and gadget
inputs before the exclusive-or replacement gives
\[
 \mathbb P_i(B=1\mid R=r,C=c,T=0)
  =\frac{(1-t_i)h_{01}+t_i h_{11}}
         {(1-t_i)(h_{00}+h_{01})+t_i(h_{10}+h_{11})}.
\]
The difference of these two fractions is
\((t_1-t_2)(h_{11}h_{00}-h_{01}h_{10})\) divided by
their positive denominator product. It is nonzero. Thus a
coordinate of the terminal response has a different conditional
law at the same full evidence assignment, while observations agree.
Restoring the moved conditioning coordinates at their action values
through the Rule 2 equalities gives different original conditional
effects. All those contexts are positive in the constructed models.

This completes the source's conditioning-marginal case division.
The active path is used through explicit independent coordinates,
so its channels may intersect the old model's vertices without
changing or exposing the unused bit. The construction avoids the
invalid general hedge inclusion and proves necessity on the
stated graph-uniform finite positive-law class.
\end{proof}
\end{annexitem}

\begin{annexitem}{annex:p16.idc-completeness}{thm:p16-idc-completeness}{Complete IDC and its conditional hedge and calculus consequences}
\begin{proof}
The exchange-closure theorem gives the same maximal moved set
\(E\) in every order. The terminal-reduction theorem equates
identification of the original conditional with identification of
the full terminal joint query. ID is complete for that graph-
uniform positive-law joint problem. Thus IDC succeeds exactly
on the identified conditional queries, and its already proved
success soundness identifies every returned formula.

The joint hedge criterion applies to the exact terminal action
and response sets, giving the stated no-hedge condition. If a
terminal hedge exists, joint nonidentification together with the
two-case conditional construction gives positive conditional
countermodels for the original query; the failed joint record
alone was not the whole argument. A successful trace consists
of Rule 2 exchanges, a successful ID derivation, marginalization
and division by a positive mass. Combining these finite derivations
proves do-calculus completeness for the stated conditional class.
No identification claim for arbitrary null conditioning contexts
or unchanged prescribed small alphabets is needed by this proof.
\end{proof}
\end{annexitem}

\begin{annexitem}{annex:p16.idc-root-graph-bridge}{thm:p16-idc-root-graph-bridge}{Active-walk conversion between two root expansions}
\begin{proof}
Let \(\Gamma\) use the model's actual latent roots. Every
latent vertex is a root, so a segment of a path between observed
vertices that passes through a latent vertex has the form
\(i\leftarrow u\to j\). Its endpoints share an actual
input and hence have a bidirected edge in \(G\). Replace
that segment by the corresponding formal-root fork in \(\widehat G\).
Conversely, for each formal-root fork choose an actual root shared
by its endpoints and replace the segment by that root's fork.
Such an actual root exists by the definition of the observed graph.

The replacement preserves arrowheads at each observed vertex and
hence its collider status. Each inserted latent occurrence is an
unconditioned noncollider, because conditioning sets are observed.
Directed ancestry from an observed collider to an observed conditioning
vertex uses only observed directed edges: a directed path from an
observed vertex cannot enter a parentless latent root. Thus collider
activation is also preserved. If several forks use the same actual
root, the result may be a walk rather than a path; the active-walk
reduction theorem \cref{thm:p16-active-walk-reduction} supplies
an active path with the same observed endpoints.

Incoming-arrow deletion at an observed action removes the same
observed directed edges and all root arrows into that vertex in
either expansion. Outgoing-arrow deletion removes its directed
arrows to observed children; root arrows point into the vertex
and are unaffected. Each surviving fork and directed edge therefore
has a corresponding surviving segment after the specified deletions.
The same conversion proves equivalence of active paths in those
modified graphs. This proves equality of the separation tests,
without replacing any actual latent input by independent pair inputs
in a probability model.
\end{proof}
\end{annexitem}

\begin{annexitem}{annex:p16.idc-rule-two-step}{thm:p16-idc-rule-two-step}{The assignments and denominators in one exchange}
\begin{proof}
The graph bridge transfers eligibility to the model's actual expanded
DAG. Apply \cref{thm:p16-do-rule-two} with baseline action
\(X=x\), proposed exchange set \(\{w\}\), retained
observations \(Z\setminus\{w\}\), and their specified values.
Its observational-side conditioning mass is exactly \(P_x(z)\).
If that is positive, the theorem supplies
\(P_{x,z_w}(z_{Z\setminus\{w\}})>0\) and the displayed
conditional equality. The tuple \(z_w\) becomes an assigned
intervention value, while every other conditioning coordinate remains
observed. The positive-context lemma supplies this support everywhere
on the main execution domain. No equality of arbitrary null
conditional versions is used.
\end{proof}
\end{annexitem}

\begin{annexitem}{annex:p16.idc-exchange-closure}{thm:p16-idc-exchange-closure}{A moral-graph proof of exchange persistence and maximality}
\begin{proof}
Work in the canonical DAG. After deleting incoming arrows at \(X\),
the action vertices are roots and are conditioned in every tested
separation. An internal path through such a root is a conditioned
noncollider, and no other vertex can use it as a conditioned
descendant. They can therefore be removed for these tests. Let
\(K=\widehat G_{\overline X}\setminus X\), put
\(H=\operatorname{an}_K(Y\cup Z)\), and let \(M\) be
the moral graph of \(K[H]\). Let \(B=H\setminus Z\),
and let \(R_Y\) be the union of the connected components of
\(M\setminus Z\) that meet \(Y\). These are graph vertex
sets and include latent vertices when they occur in \(H\).

First fix \(w\in Z\). Deleting outgoing arrows at \(w\)
does not change the ancestor set of \(Y\cup Z\). A former
ancestor path using one of those arrows can instead stop when
it reaches \(w\), itself in the target set. The converse
inclusion follows because deletion adds no paths. In the moral
graph this deletion changes only edges incident with \(w\):
its edges to children and its marriages to their other parents
disappear, while its incoming-parent edges remain. Marriages
among the other parents of any child are unchanged. Hence the
graph on \(B\) is still \(M\setminus Z\), and the only
neighbors through which \(w\) can reach it are its parents
outside \(Z\). The ancestral moral separation theorem gives
the exact criterion
\[
 w\text{ is eligible}\quad\Longleftrightarrow\quad
 \operatorname{pa}_K(w)\cap R_Y=\varnothing.
\]

Now exchange an eligible \(w\). The next reduced DAG is
\(K\setminus\{w\}\), because \(w\) is now a fixed
root that is removed from internal path tests. The new conditioning
set is \(Z'=Z\setminus\{w\}\), and its ancestor set
\(H'\) is a subset of \(H\setminus\{w\}\).
We show that the union of response-connected moral components is
still exactly \(R_Y\).

Any vertex \(t\in H\setminus(H'\cup\{w\})\) has an
old directed path to \(w\). Choose such a path. It cannot
meet another vertex of \(Y\cup Z\) before \(w\), because
that would give \(t\) a surviving path to a retained target.
Consequently it connects \(t\), in \(M\setminus Z\), to
a parent of \(w\). Eligibility puts that parent outside
\(R_Y\), so \(t\notin R_Y\). Thus no vertex of \(R_Y\)
is lost from the ancestor set.

No moral edge inside \(R_Y\) is lost either. An edge removed
because its child was \(w\) joins parents of \(w\), all
outside \(R_Y\). An edge removed because another child was
lost from \(H'\) cannot have a parent in \(R_Y\): that
parent was also adjacent to the lost child in the old moral
graph, which would put the child in \(R_Y\), a contradiction.
Observed adjacency edges between retained vertices stay in place.
The new graph is a subgraph of the old one on the retained
vertices, so it creates no new response-connected vertices. This
proves invariance of \(R_Y\).

For any remaining \(v\in Z'\), its parent set changes only
by removal of \(w\), which was already outside \(R_Y\).
Every other parent of \(v\) is in the new ancestor set since
\(v\) remains a target for ancestry. Its parent intersection
with \(R_Y\) is therefore unchanged. The criterion above
shows that its eligibility is unchanged. Repetition proves that
exactly the initially eligible vertices can be moved, in any order.

For the simultaneous statement, delete outgoing arrows from a
set \(T\subseteq Z\). The ancestor set \(H\) is again
unchanged, and the moral graph on \(B\) is still the same.
Vertices in \(T\), left as tested endpoints while \(Z\setminus T\)
is conditioned, connect to \(B\) only through parents outside
\(Z\). Edges directed between members of \(T\) have been
removed with their outgoing arrows. If every member is initially
eligible, none of those endpoint-parent connections reaches
\(R_Y\), so the simultaneous separation holds. If a member
has a parent in \(R_Y\), that parent edge and its path to
\(Y\) in \(B\) remain, so simultaneous separation fails.
Thus precisely the subsets of the initial eligible set satisfy
the simultaneous test, and its unique maximal set is \(E\).
\end{proof}
\end{annexitem}

\begin{annexitem}{annex:p16.idc-termination-soundness}{thm:p16-idc-termination-soundness}{All returns, recursive steps and the terminal normalization}
\begin{proof}
Each nonterminal IDC call moves one conditioning vertex to its
action set and strictly decreases \(|Z|\). The observed graph,
the original observed law and the response set are unchanged.
Hence there are at most the original \(|Z|\) exchanges. The
terminal ID call terminates by its already proved decreasing measure.
The empty-response return terminates immediately and is the conditional
mass of the empty tuple, namely one on every supported context.

Induct along the retained exchange sequence. The preceding Rule 2
lemma identifies each new conditional query with its parent at the
same assigned values. It also preserves support, and the positive
context lemma supplies all contexts for every model in the main
domain. At the terminal call, successful ID returns exactly the
joint intervention law \(q_a(y,w)=P_a(y,w)\) for the enlarged
action and remaining conditioning sets. Its marginal over \(y\)
is \(P_a(w)>0\). Dividing therefore gives the correct terminal
conditional, and the retained exchange equalities give the original
conditional. If ID fails, no quotient is formed and this success
proof makes no claim about that record's conditional necessity.

Every exchange is the proved Rule 2 application. The successful
terminal ID trace already supplies its finite do-calculus derivation;
marginalization and normalization are ordinary probability operations.
Composing these derivations proves the claimed successful derivation.
The first worked example below proves the final possibility concerning
an unidentified original joint law.
\end{proof}
\end{annexitem}

\begin{annexitem}{annex:p16.idc-known-denominator}{thm:p16-idc-known-denominator}{The identified-marginal case of the source necessity argument}
\begin{proof}
If the terminal joint is identified, successful ID and normalization
identify the conditional by the preceding theorem. Conversely, suppose
the original conditional is graph-uniformly identified. The Rule 2
exchange equalities then identify the terminal conditional for every
assignment: an enlarged action assignment is exactly an original
action tuple together with values of the moved conditioning coordinates.
By hypothesis the terminal conditioning marginal is also identified.
Their product identifies \(P_a(y,w)\), proving the converse.

More explicitly, if terminal joint ID fails, the joint completeness
theorem supplies a positive four-valued pair with the same observed
law and a different terminal joint at some action and response
assignment. Their conditioning marginals agree by the hypothesis
and are positive. If their conditional kernels also agreed, multiplying
by that common marginal would make the joints equal, a contradiction.
At a differing conditional assignment, the exchanged action values
can be restored as the original evidence coordinates. Rule 2 and
positivity then give two different original conditional effects.
When no conditioning remains, its marginal is the unit mass and
the same proof applies. This is the identified-denominator branch
of the source's Appendix argument, with both models and their
support kept explicit.
\end{proof}
\end{annexitem}

\begin{annexitem}{annex:p16.idc-response-rooted-failure}{thm:p16-idc-response-rooted-failure}{A response-root parity discrepancy survives conditioning}
\begin{proof}
Use the two positive binary forest models with private noise
\(E_0\) at one chosen common root and
\(0<\epsilon=\mathbb P(E_0=1)<1/2\). Extend every outside
observed coordinate independently with a fair private bit, retaining
other graph edges as ignored arguments and independent ignored
pair inputs. This is the earlier full-graph embedding construction;
its observed laws are equal and strictly positive.

Under the terminal action, the parity of the common roots is
fair in model 1 and equals \(E_0\) in model 2. All roots
lie in \(Y\), so none belongs to \(W\). The selected
root carrying \(E_0\) has no selected child, and outside
equations ignore it. Thus in model 2 the conditioning tuple
\(W\) is independent of \(E_0\), even after the action.
The conditional expected sign of root parity in model 2 is
the constant \(1-2\epsilon\) at every \(w\).

If the conditional distributions of \(Y\) given \(W\)
agreed in the two models at every stratum, the same constant
would be the conditional expected sign of root parity in model 1.
Averaging under model 1 would give its unconditional expected
sign as \(1-2\epsilon\). But its root parity is fair, giving
expected sign zero. This is a contradiction. All conditioning
strata in both models are positive by full observed positivity.
Therefore some conditional response differs. Restoring the moved
observation values through the exchange equalities supplies the
original-query discrepancy. One binary family suffices to disprove
graph-uniform identification.
\end{proof}
\end{annexitem}

\begin{annexitem}{annex:p16.idc-parent-root-gadget}{thm:p16-idc-parent-root-gadget}{A positive likelihood gadget at a conditioning root}
\begin{proof}
Take the positive forest pair with its private noise \(E_0\)
at the chosen root \(r\), of one-probability \(\epsilon\)
strictly between zero and \(1/2\). Introduce the outside
response vertex \(s\) as an independent fair bit, and add
its value modulo two to the equation at \(r\) in both models.
This uses the given allowed edge \(s\to r\). Since \(r\)
is a selected root, no other forest equation uses its value.
Conditional on either \(s\), the new forest law is the same
bijective flip of the old root coordinate in both models, so
the full observed laws still agree and are positive. Give all
other outside vertices independent positive private values and
retain the remaining declared graph edges as ignored arguments.

The action misses \(s\) and the conditioning roots. In model 1,
the parity \(T\) of all common roots equals the old fair
root parity plus \(s\). Its product sign with \(s\) therefore
has expectation zero. In model 2 the parity identity is exactly
\(T=s\oplus E_0\). Every coordinate of \(W\setminus\{r\}\)
is independent of \(s,E_0\), because these inputs affect only
the root \(r\). Since \(W\) includes all common roots,
fixing \(W=w\) fixes the parity \(T=t(w)\). Bayes' formula
with the fair prior on \(s\) consequently gives in model 2
\[
 \mathbb E_2[(-1)^s\mid W=w]
       =(1-2\epsilon)(-1)^{t(w)}.
\]
This follows directly from the two weights
\(\frac12\mathbb P(E_0=t(w))\) and
\(\frac12\mathbb P(E_0=1\oplus t(w))\); all other fixed
conditioning coordinates contribute the same factor to both.

If the conditional response laws agreed everywhere, their marginal
conditional laws on the response coordinate \(s\) would agree.
The displayed identity would then hold in model 1. Multiplying
by \((-1)^{t(w)}\) and averaging under model 1 would make
\(\mathbb E_1[(-1)^{s+T}]=1-2\epsilon\), contradicting
its zero value. Thus a conditional response differs. Strict
positive observations supply positive post-action conditioning
in both models. This is a complete positive parent-entry case
of the source's conditional countermodel construction.
\end{proof}
\end{annexitem}

\begin{annexitem}{annex:p16.idc-hedge-scope-counterexample}{thm:p16-idc-hedge-scope-counterexample}{The hedge inclusion counterexample and its full positive law}
\begin{proof}
The directed part is a path to \(W\), with at most one
child per vertex, and the bidirected part is a spanning path.
Thus \(F\) is a W-rooted C-forest. Its singleton smaller
forest \(F'\) has the same root, misses \(X\), and the
larger forest contains \(X\). The root is itself the response,
so the post-action ancestor condition holds. These are precisely
the hedge conditions. But \(L\) belongs to \(F\) and
to neither \(F'\) nor the action set, disproving the stated
general inclusion.

In the supplied models \(X=(U\oplus V)\oplus U=V\)
naturally, and the map \((U,V)\mapsto(L,X)=(U\oplus V,V)\)
is bijective. Hence \(L,X\) are independent fair bits.
In both models naturally \(W=E\), independent of those
bits. Every observed cell therefore has mass \((1/4)(3/4)=3/16\)
at \(w=0\) and \((1/4)(1/4)=1/16\) at \(w=1\).
After replacing \(X\) by \(x\), the first response is
\(x\oplus V\oplus E\), which is fair, and the second
remains \(E\). The action and every response context are
supported. The example concerns a legitimate general hedge,
so a proof needing the rejected inclusion must supply additional
hypotheses or a different construction.
\end{proof}
\end{annexitem}

\begin{annexitem}{annex:p16.idc-worked}{thm:p16-idc-worked}{The full chain and collider traces}
\begin{proof}
In the first graph, after incoming arrows at \(X\) and
outgoing arrows at \(Z\) are removed, no active path joins
\(Z\) to \(Y\) given \(X\): the only remaining route
to \(Y\) would use its deleted observed parent edge. The
canonical root bridge also verifies this in the actual root DAG.
IDC therefore exchanges \(Z\), retaining its value as an
action, and calls the unconditional subproblem with actions
\(\{X,Z\}\) and response \(\{Y\}\). ID retains
the three directed ancestors, finds the singleton Y district
already maximal in the full graph, and returns
\(p(y\mid x,z)\). The inputs at \(Y\) are private and
its only parent is \(Z\), so this equals \(p(y\mid z)\).
Its terminal denominator is one, because no conditioning remains.

For the numerical pair, naturally \(X=U,Z=E\) in both,
and \(Y=Z\oplus E_Y\). Multiplying their independent
root masses gives, for each \(x\), the four entries
\(9/32,3/32,1/32,3/32\) in the stated order. They are
positive and sum to \(1/2\) for each \(x\). Under
\(X=x\), the first model has \(Z=x\oplus U\oplus E\)
fair, whereas the second retains \(Z=E\). Thus the joint
law on \(Z,Y\) differs. Given any \(Z=z\), however,
the private \(E_Y\) remains independent in both models, giving
the displayed common conditional probabilities. This establishes
conditional success despite failure of identification for the
original joint intervention law.

For the second graph, the path \(Z\leftarrow Y\) survives
the proposed outgoing-arrow deletion at \(Z\), and is active
given \(X\). Thus \(Z\) is not eligible. The terminal
joint request has response \(\{Y,Z\}\), action \(\{X\}\).
ID retains all ancestors and splits the two post-action districts
\(\{Y\},\{Z\}\). The Y factor is its observed marginal.
For the Z factor, the full graph's containing district is
\(\{X,Z\}\). Extraction creates its kernel with \(Y\)
as a fixed parameter. The recursive graph on \(X,Z\)
has \(X\to Z,X\leftrightarrow Z\), whose whole and
remaining singleton graphs are single districts. The joint ID
call therefore fails, and IDC retains that failure without division.

Numerically, natural substitution of \(X=U\) gives
\(Z=Y\oplus E\) in both models. Independent fair \(X,Y\)
and the noise mass give \(p(x,y,z)=3/16\) at \(z=y\)
and \(1/16\) otherwise, with all eight cells positive.
Intervening on \(X\) leaves the fair \(Y=V\) unchanged.
In model 1 the extra fair \(U\) makes \(Z\) fair and
independent of \(Y\), so the conditional response-one mass
is \(1/2\) at either \(z\). In model 2 the pair remains
\(Y=V,Z=V\oplus E\); since \(Z\) is fair, Bayes'
formula gives \(1/4\) at \(z=0\) and \(3/4\) at
\(z=1\). Both conditioning masses equal \(1/2\).
These are complete positive countermodels for this failed conditional
query, not a conclusion inferred from the failed joint alone.
\end{proof}
\end{annexitem}
